\documentclass[12pt,a4paper]{article}

\usepackage[utf8]{inputenc}
\usepackage[T2A]{fontenc}
\usepackage[english,russian]{babel}
\usepackage{array}
\usepackage{booktabs}
\usepackage{geometry}
\usepackage{hyperref}
\usepackage{longtable}
\geometry{margin=25mm}
\hypersetup{unicode=true,colorlinks=true,urlcolor=blue,linkcolor=black}
\newcommand{\code}[1]{\texttt{\detokenize{#1}}}

\begin{document}

\begin{titlepage}
\begin{center}
{\large Московский авиационный институт\\
(национальный исследовательский университет)}

\vfill
Лабораторная работа по курсу «Информационный поиск»

\vspace{12mm}
{\Large \textbf{Лабораторная работа №1}}\\
{\large \textbf{Добыча корпуса документов}}

\vfill
\begin{flushright}
Студент: К.\,А. Арусланов\\
Группа: \underline{\hspace{45mm}}\\
Преподаватель: А.\,А. Кухтичев\\
Дата: \underline{\hspace{45mm}}\\
Оценка: \underline{\hspace{45mm}}
\end{flushright}

\vfill
Москва, 2026
\end{center}
\end{titlepage}

\tableofcontents
\newpage

\section{Постановка задачи}

Необходимо скачать примеры документов будущего корпуса, изучить их разметку
и метаданные, выделить основной текст, найти внешние средства поиска и
посчитать характеристики выборки. Темой корпуса выбраны русскоязычные
редакционные публикации о видеоиграх и игровой индустрии: новости, статьи,
обзоры, превью и аналитика.

\section{Источники и метод}

Кандидатами являются PlayGround.ru, StopGame.ru, IXBT.games и Игромания.
В окончательном корпусе будет не менее двух независимых источников. В
\code{urls.txt} подготовлено по две репрезентативные страницы каждого сайта.
Массовый обход на этапе лабораторной работы не выполняется.

Сбор запускается командой:

\begin{verbatim}
python main.py collect urls.txt --delay 2
\end{verbatim}

Перед запросом программа получает и сохраняет актуальный \code{robots.txt},
проверяет URL для собственного User-Agent и не запрашивает запрещённые пути.
Ответ классифицируется как документ, редирект, 404, 429, 5xx,
anti-bot/challenge или запрет robots.txt. CAPTCHA и WAF не обходятся.

Сырой HTML хранится в \code{data/gaming/raw/<source>/}, JSON с результатом
разбора --- в \code{data/gaming/parsed/<source>/}. Игровой корпус вынесен в
отдельный каталог и не смешивается с прежней темой lyrics.

\section{Единая модель документа}

Все парсеры возвращают источник, URL, заголовок, автора, дату публикации,
категорию, теги, связанные игры, основной текст и исходные метаданные. Также
записываются HTTP-статус, размеры сырого и выделенного текста, SHA-256
нормализованного текста и ошибка разбора. Неизвестные метаданные не
выдумываются. Удаляются только tracking-параметры URL, а содержательные
параметры сохраняются.

\section{Разметка и выделение текста}

На двух локально сохранённых страницах StopGame подтверждена следующая
структура:

\begin{center}
\begin{tabular}{p{0.32\linewidth}p{0.55\linewidth}}
\toprule
Данные & Разметка \\
\midrule
Основной текст & \code{article#material_content} \\
Заголовок & \code{h1}, JSON-LD \code{NewsArticle.headline} \\
Автор & JSON-LD \code{NewsArticle.author.name}, Open Graph \\
Дата & JSON-LD \code{datePublished}, \code{og:article:published_time} \\
Описание & JSON-LD \code{description}, \code{meta[name=description]} \\
Игры & \code{a.game-link[data-gamelink-game]} внутри текста \\
\bottomrule
\end{tabular}
\end{center}

Недопустимо извлекать любой тег \code{article}: отдельные такие элементы
используются для комментариев и рекомендаций. Внутри основного материала
блоки с классом, содержащим \code{_read-more}, содержат ссылки «Предыдущий
отчёт», «Подробнее по теме» и «Читай также»; они удаляются до извлечения.

Для PlayGround, IXBT.games и Игромании HTML-образцы в данной локальной среде
ещё не получены: соединение к сайтам завершалось тайм-аутом до HTTP-ответа.
Поэтому их специфические селекторы не заявляются подтверждёнными. До
сохранения 2--5 образцов каждого сайта включён только консервативный
семантический резерв: \code{[itemprop=articleBody]}, \code{article},
\code{main}. Его использование явно отмечается в JSON и должно быть заменено
после ручной проверки.

\section{Статистика проверочной выборки}

Локально сохранены две страницы StopGame: 203\,473 и 268\,147 байт HTML.
Разбор повторяется без обращения к сети:

\begin{verbatim}
python main.py parse-fixtures data/gaming/fixtures.tsv
python main.py stats
\end{verbatim}

\begin{center}
\begin{tabular}{lr}
\toprule
Показатель & Значение \\
\midrule
Количество документов & 2 \\
Успешно распарсено / ошибок & 2 / 0 \\
Размер сырого HTML & 471\,620 байт \\
Размер выделенного текста & 13\,732 байта \\
Средний raw / text & 235\,810 / 6\,866 байт \\
Медиана текста & 6\,866 байт \\
Минимум / максимум текста & 2\,144 / 11\,588 байт \\
Слов всего & 1\,131 \\
Среднее / медиана слов & 565.5 / 565.5 \\
Точных дубликатов & 0 \\
Коэффициент извлечения & 2.9117\% \\
\bottomrule
\end{tabular}
\end{center}

Это техническая выборка для проверки разметки; её нельзя экстраполировать на
миллион документов или считать финальным мульти-источниковым корпусом.

\section{Существующий поиск}

Внешний веб-поиск по ограничению \code{site:} находит документы каждого
кандидатного источника. Примеры: \code{site:playground.ru/elden\_ring/opinion
"Обзор Elden Ring"}, \code{site:stopgame.ru/newsdata "BioShock 4" "ценник
GTA VI"}, \code{site:ixbt.games/news "GTA VI" "ПК"} и
\code{site:igromania.ru/news "Второй трейлер GTA VI"}. В последнем случае
находится новость о 475 млн просмотров; в StopGame --- сохранённый материал о
Take-Two. На сайте StopGame также присутствует поле внутреннего поиска.

Недостатки внешней выдачи: сильная зависимость от свежести и популярности,
отсутствие контроля над нормализацией русских слов и опечаток, а также
невозможность отделить текст статьи от рекомендаций по одному только
сниппету.

\section{Тестирование}

Запуск тестов:

\begin{verbatim}
python -m unittest discover -s tests -v
\end{verbatim}

Выполненный прогон содержит три успешных теста: проверку селектора StopGame и
исключения связанных блоков, статистику двух fixtures без дубликатов и
канонизацию URL. После получения образцов трёх остальных источников план
дополняется теми же проверками, а также кейсами 404, 429, robots.txt и
challenge-страницы.

\section{Вывод}

Подготовлена воспроизводимая основа игрового корпуса: общий формат документа,
сохранение raw HTML, диагностика ответов, статистика и подтверждённый на
fixtures парсер StopGame. Для завершения ЛР 1 необходимо получить реальные
HTML-образцы хотя бы ещё одного источника, проверить его селекторы и собрать
малую мульти-источниковую выборку. Перед целью более одного миллиона
документов требуется отдельно посчитать допустимые URL по sitemap/RSS.

\end{document}
